% Source: https://artofproblemsolving.com/community/c3h2408513_base_number_problem
% Stable identity: aops-topic-2408513-q01
For $AB$ to be a two-digit base-$b$ numeral we need $1\le A$ and $B<b$, so $b\ge3$. Then
\[n=(b-2)b+2=b^2-2b+2=(b-1)^2+1=1\cdot(b-1)^2+0\cdot(b-1)+1.\]
The digits $1,0,1$ are allowed in base $b-1\ge2$, so the representation is $\boxed{101_{b-1}}$.

(For example, with $b=10$: $n=82$, and $82=81+1=101_9$.) Note that writing $n=(b-1)\cdot(b-1)+1$ as the two ``digits'' $b-1$ and $1$ would not be a valid base-$(b-1)$ numeral, because the largest digit in base $b-1$ is $b-2$.
